\begin{remark}[The motivic dimension formula]
\hypertarget{tag-IA11}{}\label{I:motivic-dimension-input}
\StatementRecord{I:motivic-dimension-input}{proved}{IA11}{}
\citet[Construction~3.4, p.~12, and Lemma~3.5, p.~13, arXiv v2]{bertolin2022} compute the motivic
group from its abelian point and extension data and the central
Weil-pairing bracket. For the $1$-motives used here, let $B$ be the
smallest abelian subvariety generated by these data, and let $Z(1)$
be the central torus obtained from that bracket and the remaining
algebraic lift. Their dimension formula is
\begin{equation}\hypertarget{tag-003D}{}\label{ari:eq:cp-MTdim}
 \dim G_{\mathrm{mot}}(M)
 =\dim G_{\mathrm{mot}}(A)+2\dim B+\dim Z(1).
\end{equation}
The bracket is described in their Section~6.1, pp.~26--27, and
their corrected Table~2, p.~29, with the correction on p.~28, gives unipotent dimension
three in the dependent nondeficient case. The polarization calculation below identifies these groups
for the point and extension data in this paper.
\end{remark}